Gràcies a les valuoses dades sobre les posicions dels astres que Tycho Brahe va recollir al
llarg de la seva vida, Johannes Kepler va poder formular les seves famoses tres lleis.

\begin{apartat}{1,25}
Deduïu la tercera llei de Kepler a partir de la segona llei de Newton i de la llei de
gravitació universal, suposant que els planetes descriuen moviments circulars uniformes.
\end{apartat}

\begin{apartat}{1,25}
A partir de les dades de la taula, determineu la massa del Sol.

\begin{center}
\begin{tabular}{lcc}
\textit{Planeta} & \textit{Radi de l'òrbita ($10^{9}$ m)} & \textit{Període (anys)} \\
Mercuri & 57,90 & 0,2408 \\
Venus & 108,2 & 0,6152 \\
Terra & 149,6 & 1,000 \\
Mart & 228,0 & 1,881 \\
\end{tabular}
\end{center}
\end{apartat}

\dades{$G = 6,67\times 10^{-11}\ \mathrm{N\,m^2\,kg^{-2}}$.}
